\scanpage{162}
\mspage{\unsure{1|I}}
1) \uline{Filtration du complexe de DeRham d'un schéma relatif \inserted{lisse}, et connexion de Gauss-Manin} \deleted{\unsure{en dim.}} (d'après Katz-Oda).

\diagram{\begin{tikzcd}[ampersand replacement=\&,row sep=large]
X \arrow[d,"{\substack{f\ \text{lisse [ou seulement diff. lisse,}\\\text{pour rendre exacte la}\\\text{suite (1)]}}}"]\\
S \arrow[d,"g"]\\
T
\end{tikzcd}}

On a donc une suite exacte
\[
 \deleted{\text{On a donc}}\ 0\longrightarrow f^*(\Omega^1_{S/T})\longrightarrow\Omega^1_{X/T}\longrightarrow\Omega^1_{X/S}\longrightarrow0 . \tag*{(1)}
\]
d'où une filtration $\mathrm{Fil}_{X/S}$ de l'objet gradué $\Omega^\bullet_{X/T}=\bigwedge^{*}\Omega^1_{X/T}$, avec gradué associé
\[
 \mathrm{Gr}^p_{X/S}(\Omega^\bullet_{\deleted{\unsure{S|X}}X/T})\simeq f^*(\Omega^p_{S/T})\otimes\Omega^\bullet_{X/S}[-p] \tag*{(2)}
 % Faint pencil brackets above the first factor on the left and beneath the shifted factor on the right.
\]
On constate que la filtration respecte les opérateurs \unsure{différentiel|différentiels} de $\Omega^\bullet_{X/T}$, et qu'il en est de $=$ de l'isomorphisme (2), \inserted{compte tenu que, grâce à la conn. intégrable de $f^*(\Omega^p_{S/T})$ rel : $X/S$, $f^*(\Omega^p_{S/T})\otimes\Omega^\bullet_{X/S}$ \unsure{se trouve|étant|est} bien muni d'une différentielle.} [D'ailleurs, comme $f$ est plat, on a
\[
 f^*(\Omega^p_{S/T})\simeq Lf^*(\Omega^p_{S/T})
\]
et comme $\Omega^\bullet_{X/S}[-p]$ est plat, le $\otimes$ de (2) peut se remplacer par $\overset{\mathbb L}{\otimes}$, donc on a
\[
 \mathrm{Gr}^p_{X/S}(\Omega^\bullet_{X/T})\simeq\mathbb L f^*(\Omega^p_{S/T})\overset{\mathbb L}{\otimes}\Omega^\bullet_{X/S}[-p] \,, \tag*{(2$'$)}
\]
dans la mesure où cela veut dire qqe chose].

La filtration sur $\Omega^\bullet_{X/T}$ définit $\mathbb R f_*(\Omega^\bullet_{X/T})$ comme objet d'une catégorie dérivée filtrée, (filtration encore notée $\mathrm{Fil}^\bullet_{X/S}$, d'où $\mathrm{Gr}^\bullet_{X/S}$). On
% archive sheet number: 161
